\chapterpage{Stacks and queues}{Five problems. A stack is the right answer
whenever the most recent unresolved thing is the one you need next.}

% ------------------------------------------------------------------
\problem{Valid parentheses}{Easy}

Given a string of \texttt{()}, \texttt{[]} and \texttt{\{\}}, decide whether
every bracket is closed by the matching kind, in the right order.

\idea{Push openers. On a closer, the top of the stack must be its partner. At
the end the stack must be empty.}

\solheading
\begin{lstlisting}[style=java]
public static boolean isValid(String s) {
    if (s == null) return true;
    Deque<Character> stack = new ArrayDeque<>();
    for (char c : s.toCharArray()) {
        switch (c) {
            case '(' : case '[' : case '{' : stack.push(c); break;
            case ')' : if (stack.isEmpty() || stack.pop() != '(') return false; break;
            case ']' : if (stack.isEmpty() || stack.pop() != '[') return false; break;
            case '}' : if (stack.isEmpty() || stack.pop() != '{') return false; break;
            default  : break;    // ignore anything else
        }
    }
    return stack.isEmpty();
}
\end{lstlisting}

\cost{$O(n)$ time, $O(n)$ space.}
\trap{Two endings get forgotten. Popping an empty stack on a stray closer
throws, so test \texttt{isEmpty()} first; and leaving the stack non-empty at
the end means unclosed openers, which is also invalid.}

% ------------------------------------------------------------------
\problem{Evaluate reverse Polish notation}{Medium}

Evaluate tokens like \texttt{["2","1","+","3","*"]}, which means
$(2+1)\times 3 = 9$.

\idea{Push numbers. On an operator, pop two, apply, push the result. The first
value popped is the right-hand operand.}

\solheading
\begin{lstlisting}[style=java]
public static int evalRPN(String[] tokens) {
    Deque<Integer> stack = new ArrayDeque<>();
    for (String t : tokens) {
        switch (t) {
            case "+": case "-": case "*": case "/": {
                int right = stack.pop();           // order matters
                int left  = stack.pop();
                stack.push(apply(left, right, t));
                break;
            }
            default:
                stack.push(Integer.parseInt(t));
        }
    }
    return stack.pop();
}

private static int apply(int a, int b, String op) {
    switch (op) {
        case "+": return a + b;
        case "-": return a - b;
        case "*": return a * b;
        default : return a / b;
    }
}
\end{lstlisting}

\cost{$O(n)$ time, $O(n)$ space.}
\trap{Popping the operands in the wrong order. It is invisible for
\texttt{+} and \texttt{*} and wrong for \texttt{-} and \texttt{/}, so the bug
survives your first two test cases.}

% ------------------------------------------------------------------
\problem{Next greater element}{Medium}

For every element, find the first element to its right that is larger, or
$-1$. \texttt{[2,1,2,4]} gives \texttt{[4,2,4,-1]}.

\idea{A monotonic stack of \emph{indices} still waiting for an answer. When a
new value arrives, it resolves every waiting index smaller than it. Each
index is pushed and popped once.}

\solheading
\begin{lstlisting}[style=java]
public static int[] nextGreater(int[] a) {
    int n = a.length;
    int[] out = new int[n];
    Arrays.fill(out, -1);
    Deque<Integer> waiting = new ArrayDeque<>();   // indices

    for (int i = 0; i < n; i++) {
        while (!waiting.isEmpty() && a[waiting.peek()] < a[i])
            out[waiting.pop()] = a[i];
        waiting.push(i);
    }
    return out;
}
\end{lstlisting}

\cost{$O(n)$ time --- the inner \texttt{while} is amortised, since every index
leaves the stack at most once. $O(n)$ space.}
\trap{The nested-loop version is $O(n^2)$ and looks so natural that people
write it without noticing. Also store indices, not values, or you cannot
write the answer back to the right slot.}
\alsoknown{``Daily temperatures'', where the answer is the gap \texttt{i - j}
rather than the value.}

% ------------------------------------------------------------------
\problem{Min stack}{Medium}

A stack that also reports its minimum in $O(1)$.

\idea{Keep a second stack holding the minimum \emph{as of} each push. Pushing
records the smaller of the new value and the current minimum; popping
discards both in step, so the two stacks stay aligned.}

\solheading
\begin{lstlisting}[style=java]
class MinStack {
    private final Deque<Integer> values = new ArrayDeque<>();
    private final Deque<Integer> minima = new ArrayDeque<>();

    void push(int x) {
        values.push(x);
        minima.push(minima.isEmpty() ? x : Math.min(x, minima.peek()));
    }

    int pop() {
        minima.pop();
        return values.pop();
    }

    int top()     { return values.peek(); }
    int getMin()  { return minima.peek(); }
}
\end{lstlisting}

\cost{$O(1)$ for every operation, $O(n)$ space.}
\trap{Keeping a single \texttt{int min} field instead. It is right until you
pop the minimum, and then there is no way to recover the previous one without
scanning.}

% ------------------------------------------------------------------
\problem{Simplify a Unix path}{Medium}

Turn \texttt{"/a/./b/../../c/"} into \texttt{"/c"}. Handle \texttt{.},
\texttt{..} and repeated slashes.

\idea{Split on slashes and push each real component. \texttt{..} pops,
\texttt{.} and empty components are skipped. Join what is left.}

\solheading
\begin{lstlisting}[style=java]
public static String simplify(String path) {
    Deque<String> parts = new ArrayDeque<>();     // used as a list
    for (String part : path.split("/")) {
        if (part.isEmpty() || part.equals(".")) continue;
        if (part.equals("..")) parts.pollLast();  // no-op when empty
        else parts.addLast(part);
    }
    if (parts.isEmpty()) return "/";
    StringBuilder out = new StringBuilder();
    for (String part : parts) out.append('/').append(part);
    return out.toString();
}
\end{lstlisting}

\cost{$O(n)$ time, $O(n)$ space.}
\trap{\texttt{..} at the root must do nothing, not throw ---
\texttt{pollLast} returns \texttt{null} on an empty deque where
\texttt{removeLast} throws. And the result for an empty stack is
\texttt{"/"}, not \texttt{""}.}
